\section{Годограф Найквиста }

В соответствии с вариантом 27 из \textit{Таблицы 1 инструкции по выполнению л/р} придумаем три объекта пятого
порядка $p$ полюсов передаточных функций которых вещественные, а $q$ - комплексно-
сопряженные.
\begin{itemize}
    \item Первая передаточная функция должна иметь $n$ неустойчивых полюсов у разомкнутой системы и $m$ неустойчивых полюсов у замкнутой.
    \item Вторая передаточная функция должна иметь $0$ неустойчивых полюсов у разомкнутой системы и $m$ у замкнутой.
    \item Третья передаточная функция должна иметь $n$ неустойчивых полюсов у разомкнутой системы и $0$ у замкнутой.
\end{itemize}
Для каждого объекта:
\begin{enumerate}
    \item Приведем карты (изображения размещения на комплексной координатной плоскости) нулей и полюсов разомкнутой и замкнутой систем на комплексной координатной плоскости.
    \item Выполним моделирование и приведем переходные характеристики для разомкну-
    той и замкнутой систем.
    \item Построим годограф Найквиста (АФЧХ). Найдем число оборотов годографа по ча-
    совой стрелке вокруг точки $(-1; 0)$ и через критерий Найквиста. Сравним результаты.
    \item Построим ЛАЧХ и ЛФЧХ разомкнутой системы. Определим суммарное количество переходов ЛФЧХ через критические отрезки при положительном значении
    ЛАЧХ, сравним с результатами (количеством оборотов), полученными в предыдущем пункте. Проанализируем связь устойчивости с видом ЛАФЧХ разомкнутых систем используя логарифмический критерий Найквиста.
\end{enumerate}

\begin{table}[h]
\centering
\begin{tabular}{|c|c|c|c|c|c|c|}
\hline
\textbf{Вариант} & \textbf{p} & \textbf{q} & \textbf{n} & \textbf{m} & \textbf{i} & \textbf{j} \\ \hline
27 & 1 & 4 & 4 & 2 & 13 & 3 \\ \hline
\end{tabular}
\end{table}

Придумаем 3 объекта:
Разомкнутая система имеет вид: $W(s)=\frac{R(s)}{Q(s)}$,
замкнутая - $W_3(s)=\frac{W(s)}{1+W(s)}=\frac{R(s)}{R(s)+Q(s)}=\frac{R(s)}{D(s)}$. Составим с учётом этого системы с полиномами, имеющими подходящие нам корни, найдя $R(s)$ через $Q(s)$ и $D(s)$.

\begin{itemize}
    \item 1:
Математическая модель: $4$ неустойчивых полюса у разомкнутой системы и $2$ неустойчивых полюса у замкнутой.
$$
\begin{cases}
Q(s)= (s+1)(s-1-2i)(s-1+2i)(s-2-3i)(s-2+3i),\\
Q(s)+R(s)= (s+2)(s+1-2i)(s+1+2i)(s-1-3i)(s-1+3i)
\end{cases}\Rightarrow
$$$$
\Rightarrow R(s) = 7 s^4 - 9 s^3 + 52 s^2 + 51 s + 35
$$$$
W(s)=\frac{7 s^4 - 9 s^3 + 52 s^2 + 51 s + 35}{s^5 - 5 s^4 + 20 s^3 - 20 s^2 + 19 s + 65}
$$$$
W_3(s)=\frac{7 s^4 - 9 s^3 + 52 s^2 + 51 s + 35}{s^5 + 2 s^4 + 11 s^3 + 32 s^2 + 70 s + 100}
$$


\item 2:
Математическая модель:
Математическая модель: $0$ неустойчивых полюсов у разомкнутой системы и $2$ неустойчивых полюса у замкнутой.
$$
\begin{cases}
    Q(s)=(s+1)(s+1-i)(s+1+i)(s+2-i)(s+2+i) \\
    Q(s)+R(s)=(s+2)(s+2-3i)(s+2+3i)(s-4+i)(s-4-i)
\end{cases}
$$$$
\Rightarrow R(s) = -9 s^4 - 31 s^3 - 73 s^2 + 121 s + 432
$$$$
W(s)=\frac{-9 s^4 - 31 s^3 - 73 s^2 + 121 s + 432}{s^5 + 7 s^4 + 21 s^3 + 33 s^2 + 28 s + 10}
$$$$
W_3(s)=\frac{-9 s^4 - 31 s^3 - 73 s^2 + 121 s + 432}{s^5 - 2 s^4 - 10 s^3 - 40 s^2 + 149 s + 442}
$$


\item 3:
Математическая модель:
Математическая модель: $4$ неустойчивых полюса у разомкнутой системы и $0$ неустойчивых полюсов у замкнутой.
$$
\begin{cases}
    Q(s)=(s+1)(s-1-i)(s-1+i)(s-2-i)(s-2+i) \\
    Q(s)+R(s)=(s+2)(s+2-4i)(s+2+4i)(s+3-i)(s+3+i)
\end{cases}
$$$$
\Rightarrow R(s) = 17 s^4 + 65 s^3 + 271 s^2 + 528 s + 390
$$$$
W(s)=\frac{17 s^4 + 65 s^3 + 271 s^2 + 528 s + 390}{s^5 - 5 s^4 + 9 s^3 - 3 s^2 - 8 s + 10}
$$$$
W_3(s)=\frac{17 s^4 + 65 s^3 + 271 s^2 + 528 s + 390}{s^5 + 12 s^4 + 74 s^3 + 268 s^2 + 520 s + 400}
$$


\end{itemize}


\textbf{Построим карты для разомкнутой и замкнутой систем (см. рис. \ref{c1}-\ref{c3}).}

\begin{figure}[H]
    \centering
    \includegraphics[width=1\linewidth]{lab_6/images/pz_combined1.png}
    \caption{Карта (объект 1)}
    \label{c1}
\end{figure}
\begin{figure}[H]
    \centering
    \includegraphics[width=1\linewidth]{lab_6/images/pz_combined2.png}
    \caption{Карта (объект 2)}
    \label{c2}
\end{figure}
\begin{figure}[H]
    \centering
    \includegraphics[width=1\linewidth]{lab_6/images/pz_combined3.png}
    \caption{Карта (объект 3)}
    \label{c3}
\end{figure}

Проведем моделирование.\\

\textbf{Переходная характеристика для разомкнутой и замкнутой систем (см. рис. \ref{s1}-\ref{s3}).}
\begin{figure}[H]
    \centering
    \includegraphics[width=0.62\linewidth]{lab_6/images/step_combined1.png}
    \caption{Переходная характеристика (объект 1)}
    \label{s1}
\end{figure}
\begin{figure}[H]
    \centering
    \includegraphics[width=0.62\linewidth]{lab_6/images/step_combined2.png}
    \caption{Переходная характеристика (объект 2)}
    \label{s2}
\end{figure}
\begin{figure}[H]
    \centering
    \includegraphics[width=0.62\linewidth]{lab_6/images/step_combined3.png}
    \caption{Переходная характеристика (объект 3)}
    \label{s3}
\end{figure}

\textbf{Годограф Найквиста (см. рис. \ref{n1}-\ref{n3}).}
\begin{figure}[H]
    \centering
    \includegraphics[width=0.6\linewidth]{lab_6/images/nyquist1.png}
    \caption{Годограф Найквиста (объект 1)}
    \label{n1}
\end{figure}
Число оборотов вокруг точки (против ч.с.) (-1, 0): 2\\
Критерий Найквиста:$m = n + N$,
где\\
1. $m$ –число неустойчивых полюсов замкнутой сисемы\\
2. $n$ – Число неустойчивых полюсов разомкнутой системы\\
3. $N$ –Число оборотов годографа по часовой стрелке вокруг точки (-1,0)\\
Сравнение с результатами по критерию Найквиста: $2=4-2$ - сходится.

\begin{figure}[H]
    \centering
    \includegraphics[width=0.6\linewidth]{lab_6/images/nyquist2.png}
    \caption{Годограф Найквиста (объект 2)}
    \label{n2}
\end{figure}
Число оборотов вокруг точки (против ч.с.) (-1, 0): -2\\
Сравнение с результатами по критерию Найквиста: $2=0+2$ - сходится.

\begin{figure}[H]
    \centering
    \includegraphics[width=0.6\linewidth]{lab_6/images/nyquist3.png}
    \caption{Годограф Найквиста (объект 3)}
    \label{n3}
\end{figure}
Число оборотов вокруг точки (против ч.с.)
(-1, 0): 4\\
Сравнение с результатами по критерию Найквиста: $0=4-4$ - сходится.


\textbf{ЛАЧХ и ЛФЧХ разомкнутой системы (см. рис. \ref{l1}-\ref{l3}).}

\begin{figure}[H]
    \centering
    \includegraphics[width=0.9\linewidth]{lab_6/images/bode1.png}
    \caption{ЛАЧХ и ЛФЧХ (объект 1)}
    \label{l1}
\end{figure}
Суммарное количество переходов ЛФЧХ через критические отрезки при положительном значении ЛАЧХ: 1\\
\textbf{Анализ связи устойчивости с видом ЛАФЧХ разомкнутых систем используя логарифмический критерий Найквиста:}
Для разомкнутой системы число правых корней характеристического уравнения
$r = 4$, следовательно, согласно логарифмическому критерию Найквиста должно
выполняться условие $N_+ - N_- = r/2 = 2$. По графику ЛФЧХ данное условие не выполняется, что подтверждает неустойчивость замкнутой системы.

\textbf{Сравнение с результатами по годографу Найквиста:}\\
Число оборотов годографа по часовой стрелке: \( N = m - n = 2 - 4 = -2 \) (т.е. 2 оборота против часовой стрелки).
Наблюдается связь: \( N_+ - N_- = -\frac{N}{2} \), что подтверждается: \( 1 = -\frac{-2}{2} = 1 \).
Оба метода дают согласованный результат.

\textbf{Анализ изменения количества неустойчивых полюсов:}
\begin{itemize}
    \item Разомкнутая система: $n = 4$ неустойчивых полюса
    \item Замкнутая система: $m = 2$ неустойчивых полюса
    \item Изменение: $\Delta = m - n = -2$ (уменьшилось на 2)
    \item Годограф делает $N = -2$ оборота (2 оборота против часовой стрелки)
    \item Проверка по критерию Найквиста: $m = n + N = 4 + (-2) = 2$ - выполняется
\end{itemize}

Замыкание системы уменьшило количество неустойчивых полюсов с 4 до 2, 
но система осталась неустойчивой. Годограф Найквиста, делая 2 оборота против часовой стрелки 
вокруг точки $(-1; 0)$, указывает на уменьшение количества неустойчивых полюсов на 2.


\begin{figure}[H]
    \centering
    \includegraphics[width=0.9\linewidth]{lab_6/images/bode2.png}
    \caption{ЛАЧХ и ЛФЧХ (объект 2)}
    \label{l2}
\end{figure}
Суммарное количество переходов ЛФЧХ через критические отрезки при положительном значении ЛАЧХ: -1\\
\textbf{Анализ связи устойчивости с видом ЛАФЧХ разомкнутых систем используя логарифмический критерий Найквиста:}
Для разомкнутой системы число правых корней
характеристического уравнения равно $r = 0$, следовательно, согласно логарифмическому критерию Найквиста должно
выполняться условие $N_+ - N_- = 0$. По ЛФЧХ при положительном значении ЛАЧХ наблюдается $-1/2$ переходов через критические отрезки, что подтверждает наличие двух неустойчивых полюсов у замкнутой системы.

\textbf{Сравнение с результатами по годографу Найквиста:}\\
Число оборотов годографа по часовой стрелке: \( N = m - n = 2 - 0 = 2 \) (т.е. 2 оборота по часовой стрелке).
Наблюдается связь: \( N_+ - N_- = -\frac{N}{2} \), что подтверждается: \( -1 = -\frac{2}{2} = -1 \).
Оба метода дают согласованный результат.


\textbf{Анализ изменения количества неустойчивых полюсов:}
\begin{itemize}
    \item Разомкнутая система: $n = 0$ неустойчивых полюсов (система устойчива)
    \item Замкнутая система: $m = 2$ неустойчивых полюса
    \item Изменение: $\Delta = m - n = +2$ (появилось 2 неустойчивых полюса)
    \item Годограф делает $N = +2$ оборота (2 оборота по часовой стрелке)
    \item Проверка по критерию Найквиста: $m = n + N = 0 + 2 = 2$ - выполняется
\end{itemize}

Замыкание устойчивой разомкнутой системы привело к потере устойчивости 
и появлению 2 неустойчивых полюсов. Годограф Найквиста, делая 2 оборота по часовой стрелке, 
показывает, что при замыкании система становится неустойчивой.


\begin{figure}[H]
    \centering
    \includegraphics[width=0.9\linewidth]{lab_6/images/bode3.png}
    \caption{ЛАЧХ и ЛФЧХ (объект 3)}
    \label{l3}
\end{figure}
Суммарное количество переходов ЛФЧХ через критические отрезки при положительном значении ЛАЧХ: 2\\
\textbf{Анализ связи устойчивости с видом ЛАФЧХ разомкнутых систем используя логарифмический критерий Найквиста:}
Для разомкнутой системы число правых корней равно $r = 4$, следовательно, должно выполняться равенство $N_+ - N_- = r/2 = 2$. Это ранво числу переходов ЛФЧХ через критические отрезки и подтверждает устойчивость замкнутой системы.

\textbf{Сравнение с результатами по годографу Найквиста:}\\
Число оборотов годографа по часовой стрелке: \( N = m - n = 0 - 4 = -4 \) (т.е. 4 оборота против часовой стрелки).
Наблюдается связь: \( N_+ - N_- = -\frac{N}{2} \), что подтверждается: \( 2 = -\frac{-4}{2} = 2 \).
Оба метода дают согласованный результат.

\textbf{Анализ изменения количества неустойчивых полюсов:}

\begin{itemize}
    \item Разомкнутая система: $n = 4$ неустойчивых полюса
    \item Замкнутая система: $m = 0$ неустойчивых полюсов (система устойчива)
    \item Изменение: $\Delta = m - n = -4$ (исчезли 4 неустойчивых полюса)
    \item Годограф делает $N = -4$ оборота (4 оборота против часовой стрелке)
    \item Проверка по критерию Найквиста: $m = n + N = 4 + (-4) = 0$ - выполняется
\end{itemize}

Замыкание неустойчивой разомкнутой системы стабилизировало её - все 
4 неустойчивых полюса перешли в левую полуплоскость. Годограф Найквиста, делая 
4 оборота против часовой стрелки, подтверждает полную стабилизацию системы.



\section{Коэффициент усиления}

В соответствии с вариантом 27 из \textit{Таблицы 1 инструкции по выполнению л/р} возьмем значение $i$ и соответствующие ему передаточные функции $W_1(s)$ и $W_2(s)$ (Таблица 2). Добавим к каждой
функции коэффициент усиления $k$. Считаем $k$ положительным. Для полученных систем:
\begin{enumerate}
    \item Построим годограф Найквиста для значения коэффициента усиления $k = 1$.
    \item Рассмотрим, как влияет коэффициент усиления $k$ на кривую годографа, взяв не
    менее трех значений $k$.
    \item Найдем зависимость количества неустойчивых полюсов замкнутой системы относительно значений коэффициента усиления $k$. Найдем пределы значений коэффициента усиления $k$ относительно которых замкнутая система устойчива.
    \item Определим значение запаса устойчивости по амплитуде (если применимо).
    \item Выполним моделирование и приведем переходные характеристики замкнутой системы при значениях коэффициента $k$ для каждой из определенных областей значений, соответствующих как устойчивой системе, так и неустойчивой.
\end{enumerate}

\begin{table}[h]
\centering
\begin{tabular}{|c|c|c|}
\hline
\textbf{i} & $\mathbf{W_1(s)}$ & $\mathbf{W_2(s)}$ \\ \hline
13 & $\displaystyle \frac{s-1}{s^2+3s+1}$ & $\displaystyle \frac{10s^3-3s^2+13s-2}{10s^3+8s^2-5s+4}$ \\ \hline
\end{tabular}
\end{table}

С коэффициентом $k$ разомкнутая $i-$тая ПФ будет иметь вид: $W_\text{раз i}(s)=k\cdot W_i(s)=k\frac{R(s)}{Q(s)}$, замкнутая:  
$W_\text{з i}(s)=\frac{k\cdot W_i(s)}{1+k\cdot W_i(s)}=\frac{kR(s)}{kR(s)+Q(s)}$.\\

\subsection{Первая система}
$$W_\text{раз}(s)=\frac{ks-k}{s^2+3s+1}=
\frac{k(s-1)}{\left(s-\frac{\sqrt{5}-3}{2}\right)\left(s+\frac{3+\sqrt{5}}{2}\right)}$$
$$W_\text{з}(s)=\frac{ks-k}{s^2+(3+k)s+1-k}=
\frac{k(s-1)}{\left(s+\frac{(3+k)-\sqrt{(3+k)^2-4(1-k)}}{2}\right)\left(s+\frac{(3+k)+\sqrt{(3+k)^2-4(1-k)}}{2}\right)}$$

Замкнутая система будет устойчива при 
$$
\begin{cases}
(3+k)-\sqrt{(3+k)^2-4(1-k)}>0 \\
(3+k)+\sqrt{(3+k)^2-4(1-k)}>0
\end{cases}\Rightarrow
\begin{cases}
(3+k)^2>(3+k)^2-4(1-k) \\
(3+k)>-\sqrt{(3+k)^2-4(1-k)}
\end{cases}\Rightarrow
k<1
$$

Построим годограф Найквиста (см. рис. \ref{2n}-\ref{2n2}).
\begin{figure}[H]
    \centering
    \includegraphics[width=0.5\linewidth]{lab_6/images2/nyquist_W1_k_0.7.png}
    \caption{Годограф Найквиста $W_1(s)$}
    \label{2n}
\end{figure}
\begin{figure}[H]
    \centering
    \includegraphics[width=0.5\linewidth]{lab_6/images2/nyquist_W1_k_1.0.png}
    \caption{Годограф Найквиста $W_1(s)$}
    \label{2n2r1}
\end{figure}
\begin{figure}[H]
    \centering
    \includegraphics[width=0.5\linewidth]{lab_6/images2/nyquist_W1_k_1.5.png}
    \caption{Годограф Найквиста $W_1(s)$}
    \label{2n1}
\end{figure}
\begin{figure}[H]
    \centering
    \includegraphics[width=0.5\linewidth]{lab_6/images2/nyquist_W1_k_3.0.png}
    \caption{Годограф Найквиста $W_1(s)$}
    \label{2n2}
\end{figure}
Видим, что при увеличении коэффициента $k$ кривая годографа увеличивается в масштабе. Количество неустойчивых полюсов замкнутой системы увеличивается при пересечении кривой годографа точки $(-1, 0)$ (Важное уточнение - в том случае, когда она закручивается по часовой стрелке, как в нашем случае). При достижении критического значения $k=1$ у нашей замкнутой системы появляется неустойчивый полюс. 

Значение запаса устойчивости по амплитуде равно критическому допустимому коэффициенту $A_3=k=1$.

Переходные характеристики для коэффициентов $k$, соответствующих различным областям значений коэффициента представлены на рис. \ref{pf21}.
\begin{figure}[H]
    \centering
    \includegraphics[width=0.9\linewidth]{lab_6/images2/step_W1.png}
    \caption{Переходные характеристики $W_1(s)$}
    \label{pf21}
\end{figure}
Видим, что моделирование подтверждает теоретические ожидания и при увеличении коэффициента $k$ система становится менее устойчивой. Перестает быть устойчивой она при $k=1$.\\
\textbf{Зависимость количества неустойчивых полюсов от $k$:}
\begin{itemize}
    \item При $k < 1: m = 0$ (система устойчива)
    \item При $k = 1: m = 1$ (система неустойчива)
    \item При $k > 1: m = 2$ (система неустойчива)
\end{itemize}

\subsection{Вторая система}
$$W_\text{раз}(s)=\frac{k(10s^3-3s^2+13s-2)}{10s^3+8s^2-5s+4}$$
$$W_\text{з}(s)=\frac{k(10s^3-3s^2+13s-2)}{(10+10k)s^3+(8-3k)s^2+(-5+13k)s+(4-2k)}$$

$$D(s)=(10+10k)s^3+(8-3k)s^2+(-5+13k)s+(4-2k)$$

Воспользуемся критерием Гурвица.\\

Матрица Гурвица:
$$A= \begin{pmatrix}
8-3k & 4-2k & 0 \\
10+10k & -5+13k & 0 \\
0 & 8-3k & 4-2k
\end{pmatrix}$$
$$
10+10k>0 \Rightarrow k>-1
$$$$
8-3k>0 \Rightarrow k<\frac{8}{3}
$$$$
-5+13k>0 \Rightarrow k>\frac{5}{13}
$$$$
4-2k>0 \Rightarrow k<2
$$$$
\Delta_2 = (8-3k)(-5+13k) - (10+10k)(4-2k)>0
$$

$$(8-3k)(-5+13k) - (10+10k)(4-2k) = -40+104k+15k-39k^2 - (40-20k+40k-20k^2)$$
$$= -40+119k-39k^2 - (40+20k-20k^2) = -40+119k-39k^2 -40-20k+20k^2$$
$$=-80+99k-19k^2$$

$$k_{1,2}=\frac{-99\pm\sqrt{99^2-4\cdot(-19)\cdot(-80)}}{2\cdot(-19)}=\frac{-99\pm\sqrt{9801-6080}}{-38}=\frac{-99\pm\sqrt{3721}}{-38}$$

$$k_1=\frac{-99+61}{-38}=\frac{-38}{-38}=1,\quad k_2=\frac{-99-61}{-38}=\frac{-160}{-38}=4.2105$$

Условие $-19k^2+99k-80>0$ выполняется при $1<k<4.2105$.

$$
\begin{cases}
k>-1\\
k<\frac{8}{3}\\
k>\frac{5}{13}\\
k<2\\
1<k<4.2105
\end{cases}
$$

Замкнутая система будет устойчива при $1<k<2$.


Построим годограф Найквиста (см. рис. \ref{2n5}-\ref{2n22}).
\begin{figure}[H]
    \centering
    \includegraphics[width=0.5\linewidth]{lab_6/images2/nyquist_W2_k_0.5.png}
    \caption{Годограф Найквиста $W_2(s)$}
    \label{2n5}
\end{figure}
\begin{figure}[H]
    \centering
    \includegraphics[width=0.5\linewidth]{lab_6/images2/nyquist_W2_k_1.0.png}
    \caption{Годограф Найквиста $W_2(s)$}
    \label{2n2r12}
\end{figure}
\begin{figure}[H]
    \centering
    \includegraphics[width=0.5\linewidth]{lab_6/images2/nyquist_W2_k_1.5.png}
    \caption{Годограф Найквиста $W_2(s)$}
    \label{2n12}
\end{figure}
\begin{figure}[H]
    \centering
    \includegraphics[width=0.5\linewidth]{lab_6/images2/nyquist_W2_k_2.0.png}
    \caption{Годограф Найквиста $W_2(s)$}
    \label{2nf12}
\end{figure}
\begin{figure}[H]
    \centering
    \includegraphics[width=0.5\linewidth]{lab_6/images2/nyquist_W2_k_3.0.png}
    \caption{Годограф Найквиста $W_2(s)$}
    \label{2n22}
\end{figure}
Видим, что годограф снова растягивается при увеличении $k$. При значении $1<k<2$ система устойчива. Графики годографов подтверждают теорию - если провести линию от точки $(-1,0)$ и посчитать количество пересечений годографа с учетом знака, можно сделать вывод о том стала замкнутая система устойчивее, или наоборот. Неустойчивые полюса появляются при выходе $k$ из области $1<k<2$.

\textbf{Зависимость количества неустойчивых полюсов от $k$:}
\begin{itemize}
    \item При $k < 1: m = 2$ (система неустойчива)
    \item При $k = 1: m = 0$ (граница устойчивости)
    \item При $1<k <2: m = 0$ (система устойчива)
    \item При $k = 2: m = 1$ (система неустойчива)
    \item При $k > 2: 1\leq m \leq3 $ (система неустойчива)
\end{itemize}

Переходные характеристики для коэффициентов $k$, соответствующих различным областям значений коэффициента представлены на рис. \ref{pf251}.

Видим, что моделирование подтверждает теоретические ожидания и при попадании $k$ в область $1<k<2$ система становится устойчивой. Перестает быть устойчивой при выходе коэффициента из этой зоны.

\begin{figure}[H]
    \centering
    \includegraphics[width=1\linewidth]{lab_6/images2/step_W2.png}
    \caption{Переходные характеристики $W_2(s)$}
    \label{pf251}
\end{figure}

Исходная разомкнутая система $W_2(s)$ с $k=1$ имеет полюса, которые при замыкании ($\tau=0$) делают замкнутую систему неустойчивой. Сама устойчивость замкнутой системы достигается только в интервале $1<k<2$, при этом система теряет устойчивость как при уменьшении $k$ ниже 1, так и при увеличении $k$ выше 2. Понятие запаса по амплитуде (как минимального изменения коэффициента усиления, приводящего к потере устойчивости) не применимо.

\section{Запаздывание}

В соответствии с вариантом 27 из \textit{Таблицы 1 инструкции по выполнению л/р} возьмем значение $j$ и соответствующие ему передаточные функции $W_3(s)$ и $W_4(s)$ (таблица 3). Добавим к каждой
функции звено чистого запаздывания $e^{-\tau s}$. Для полученных систем:
\begin{enumerate}
    \item Построим годограф Найквиста для значений коэффициента запаздывания $\tau = 0$
    и $\tau = 0.5$.
    \item Рассмотрим, как влияет коэффициент запаздывания $\tau$ на кривую годографа, взяв
    не менее трех значений $\tau$.
    \item Найдем границы значений коэффициента запаздывания $\tau$ относительно которых
    замкнутая система устойчива.
    \item Определим значение запаса устойчивости по фазе (если применимо).
    \item Выполним моделирование и приведем переходные характеристики замкнутой си-
    стемы при значениях коэффициента $\tau$ для каждой из определенных областей значений, соответствующих как устойчивой системе, так и неустойчивой.
\end{enumerate}

\begin{table}[h]
\centering
\begin{tabular}{|c|c|c|}
\hline
\textbf{j} & $\mathbf{W_3(s)}$ & $\mathbf{W_4(s)}$ \\ \hline
3 & $\displaystyle \frac{2s+9}{s^2+s+9}$ & $\displaystyle \frac{10s^2-10s+13}{10s^3+37s+25}$ \\ \hline
\end{tabular}
\end{table}

С учетом звена чистого запаздывания разомкнутая передаточная функция имеет вид $W_{\text{раз }i}(s)=W_i(s)e^{-\tau s}$
а замкнутая: $W_{\text{з }i}(s)=\frac{W_i(s)e^{-\tau s}}{1+W_i(s)e^{-\tau s}}=\frac{e^{-\tau s}R(s)}{e^{-\tau s}R(s)+Q(s)}$.


\subsection{Первая система}

$$W_{\text{раз}}(s)=\frac{2s+9}{s^2+s+9}e^{-\tau s}$$
$$W_{\text{з}}(s)=\frac{(2s+9)e^{-\tau s}}{s^2+s+9+(2s+9)e^{-\tau s}}$$

$$W_{\text{раз 3}}(i\omega) = \frac{2i\omega+9}{-\omega^2+i\omega+9} = \frac{9 + i2\omega}{(9-\omega^2) + i\omega}=\frac{9 + i2\omega}{(9-\omega^2) + i\omega} \cdot \frac{(9-\omega^2) - i\omega}{(9-\omega^2) - i\omega}
$$

Числитель:
$$(9 + i2\omega) \cdot [(9-\omega^2) - i\omega] = 9(9-\omega^2) - 9i\omega + i2\omega(9-\omega^2) - i^2 2\omega^2$$
$$= (81 - 9\omega^2) - 9i\omega + i18\omega - i2\omega^3 + 2\omega^2$$
$$= (81 - 7\omega^2) + i(9\omega - 2\omega^3)$$

Знаменатель:
$$[(9-\omega^2) + i\omega] \cdot [(9-\omega^2) - i\omega] = (9-\omega^2)^2 - (i\omega)^2 = (9-\omega^2)^2 + \omega^2 = \omega^4 - 17\omega^2 + 81$$

$$
W_{\text{раз 3}}(i\omega)=\frac{(81 - 7\omega^2) + i(9\omega - 2\omega^3)}{\omega^4 - 17\omega^2 + 81}$$

$$U(\omega) = \frac{81 - 7\omega^2}{\omega^4 - 17\omega^2 + 81},\ V(\omega) = \frac{9\omega - 2\omega^3}{\omega^4 - 17\omega^2 + 81}$$

Амплитудная характеристика:
$$A_3(\omega) = \sqrt{U^2(\omega) + V^2(\omega)} = \sqrt{\frac{(81-7\omega^2)^2 + (9\omega-2\omega^3)^2}{(\omega^4 - 17\omega^2 + 81)^2}}$$

Найдём частоту среза $\omega_{\text{ср}}$ из условия $A_3(\omega_{\text{ср}}) = 1$:
$$(81-7\omega^2)^2 + (9\omega-2\omega^3)^2 = (\omega^4 - 17\omega^2 + 81)^2$$
$$\omega^8 - 34\omega^6 + 451\omega^4 - 2754\omega^2 + 6561 - 4\omega^6 - 13\omega^4 + 1053\omega^2 - 6561=0$$
$$ \omega^2(\omega^6 - 38\omega^4 + 438\omega^2 - 1701)=0$$
$$ (\omega^6 - 38\omega^4 + 438\omega^2 - 1701)=0$$

Корни данного уравнения: $\omega_{1,2}=\pm\sqrt{21}$, нам подходит $\omega=\sqrt{21}$.
 
$$\varphi_\text{з}=\pi + \varphi(\sqrt{21})=\pi +\operatorname{atan2}\left(\frac{9\sqrt{21} - 42\sqrt{21}}{21^2 - 17\cdot 21 + 81},  \frac{81 - 7\cdot 21}{21^2 - 17\cdot 21 + 81}\right)\approx\pi-1.98\approx1.162$$
$$\tau=\frac{\varphi_\text{з}}{\omega_\text{ср}}=\frac{1.162}{\sqrt{21}}\approx0.25357$$

При значении $\tau\approx0.25357$ замкнутая система теряет устойчивость.\\
Запас устойчивости по фазе равен $\varphi_\text{з}\approx1.162$.

Для системы без запаздывания ($\tau = 0$) запас по фазе $\varphi_\text{з}$ определяется как угол, на который можно повернуть (добавить отрицательного фазового сдвига) к фазовой характеристике на частоте среза $\omega_\text{ср}$, прежде чем система достигнет границы устойчивости ($-180^\circ$). Из расчетов выше, критический фазовый сдвиг, вносимый запаздыванием, равен $\varphi_\text{кр} \approx 1.162\ \text{рад} \approx 66.6^\circ$. Следовательно, для исходной системы ($\tau=0$) запас устойчивости по фазе составляет $\boldsymbol{\approx 66.6^\circ}$. Это означает, что система сохраняет устойчивость при введении дополнительного фазового запаздывания, не превышающего этой величины.


Построим годограф Найквиста (см. рис. \ref{5n}-\ref{7n}).
\\
Как видим, увеличение $\tau$ раскручивает график годографа Найквиста, при достижении критического значения коэффициента, график переходит через точку $(-1,0)$, оставляя её внутри графика, и делает систему неустойчивой.
\begin{figure}[H]
    \centering
    \includegraphics[width=0.6\linewidth]{lab_6/images3/nyquist_W3_tau_0.00.png}
    \caption{Годограф Найквиста $W_3(s)$}
    \label{5n}
\end{figure}
\begin{figure}[H]
    \centering
    \includegraphics[width=0.6\linewidth]{lab_6/images3/nyquist_W3_tau_0.25.png}
    \caption{Годограф Найквиста $W_3(s)$}
    \label{6n}
\end{figure}
\begin{figure}[H]
    \centering
    \includegraphics[width=0.6\linewidth]{lab_6/images3/nyquist_W3_tau_0.50.png}
    \caption{Годограф Найквиста $W_3(s)$}
    \label{7n}
\end{figure}


Переходные характеристики для коэффициентов $\tau$, соответствующих различным областям значений коэффициента представлены на рис. \ref{t1}.

Моделирование также подтверждает теоретические расчёты, видим, что при значении $\tau$, близком к граничному, (0.25) в системе уже появляются сильные колебания в течение переходного процесса, а при коэффициенте большем критического значения устойчивость теряется.

\begin{figure}[H]
    \centering
    \includegraphics[width=0.8\linewidth]{lab_6/images3/step_W3.png}
    \caption{Переходные характеристики $W_3(s)$}
    \label{t1}
\end{figure}


\subsection{Вторая система}

$$W_{\text{раз}}(s)=\frac{10s^2-10s+13}{10s^3+37s+25}e^{-\tau s}$$
$$W_{\text{з}}(s)=\frac{(10s^2-10s+13)e^{-\tau s}}{10s^3+37s+25+(10s^2-10s+13)e^{-\tau s}}$$

$$W_{\text{раз 4}}(i\omega) = \frac{10(i\omega)^2 - 10i\omega + 13}{10(i\omega)^3 + 37i\omega + 25} = \frac{-10\omega^2 + 13 - i10\omega}{25 + i(37\omega - 10\omega^3)}$$

$$W_{\text{раз 4}}(i\omega)=\frac{-10\omega^2 + 13 - i10\omega}{25 + i(37\omega - 10\omega^3)} \cdot \frac{25 - i(37\omega - 10\omega^3)}{25 - i(37\omega - 10\omega^3)}=\frac{(-10\omega^2 + 13 - i10\omega)(25 - i(37\omega - 10\omega^3))}{625 + (37\omega - 10\omega^3)^2}
$$
Числитель:
$$(-10\omega^2 + 13 - i10\omega) \cdot [25 - i(37\omega - 10\omega^3)] =$$
$$= (-10\omega^2 + 13)\cdot 25 + (-10\omega^2 + 13)\cdot[-i(37\omega - 10\omega^3)] - i10\omega\cdot 25 - i10\omega\cdot[-i(37\omega - 10\omega^3)]$$
$$= (-250\omega^2 + 325) - i(-10\omega^2 + 13)(37\omega - 10\omega^3) - i250\omega + 10\omega(37\omega - 10\omega^3)$$
$$= (-250\omega^2 + 325 + 370\omega^2 - 100\omega^4) + i[-\cdot(-370\omega^3 + 100\omega^5 + 481\omega - 130\omega^3) - 250\omega]$$
$$= (325 + 120\omega^2 - 100\omega^4) + i[-370\omega^3 + 100\omega^5 + 481\omega - 130\omega^3 - 250\omega]$$
$$= (325 + 120\omega^2 - 100\omega^4) + i(100\omega^5 - 500\omega^3 + 231\omega)$$

$$
W_{\text{раз }4}(i\omega)=\frac{(325 + 120\omega^2 - 100\omega^4) + i(100\omega^5 - 500\omega^3 + 231\omega)}{625 + (37\omega - 10\omega^3)^2}$$

$$U(\omega) = \frac{325 + 120\omega^2 - 100\omega^4}{625 + (37\omega - 10\omega^3)^2},\quad V(\omega) = \frac{100\omega^5 - 500\omega^3 + 231\omega}{625 + (37\omega - 10\omega^3)^2}$$

Амплитудная характеристика:
$$A_4(\omega) = \sqrt{U^2(\omega) + V^2(\omega)} = \sqrt{\frac{(325 + 120\omega^2 - 100\omega^4)^2 + (100\omega^5 - 500\omega^3 + 231\omega)^2}{[625 + (37\omega - 10\omega^3)^2]^2}}$$

Найдём частоту среза $\omega_{\text{ср}}$ из условия $A_3(\omega_{\text{ср}}) = 1$:
$$(325 + 120\omega^2 - 100\omega^4)^2 + (100\omega^5 - 500\omega^3 + 231\omega)^2 = [625 + (37\omega - 10\omega^3)^2]^2$$

Корни данного уравнения, которые нам подходят: $\omega\approx1.7975, \omega\approx2.33$.

$$\varphi_\text{з1}=\pi + \varphi(1.7975)
\approx\pi-2.06675\approx1.0748$$
$$\varphi_\text{з2}=\pi + \varphi(2.33)
\approx\pi-2.634\approx0.5076$$

$$\tau_1=\frac{\varphi_\text{з1}}{\omega_\text{ср1}}=\frac{1.0748}{1.7975}\approx0.598$$
$$\tau_2=\frac{\varphi_\text{з2}}{\omega_\text{ср2}}=\frac{0.5076}{2.33}\approx0.2178$$
При значении $0.2178<\tau<0.598$ замкнутая система будет устойчивой.
\\
В данном случае разомкнутая система $W_4(s)$ имеет полюса с положительной вещественной частью, поэтому она неустойчива.
При $\tau = 0$ замкнутая система также неустойчива, так как $\tau=0$ не входит в интервал устойчивости $0.218 < \tau < 0.598$. Следовательно, понятие запаса устойчивости по фазе (как дополнительного фазового сдвига, который можно добавить до потери устойчивости) для данной системы при $\tau=0$ не применимо.

Построим годограф Найквиста (см. рис. \ref{5yn}-\ref{7yn}).
\\
Результаты снова подтверждают теорию, только на этот раз когда точка оказывается внутри графика годографа точки $(-1,0)$, система становится устойчивой, так как линии годографа закручиваются против часовой стрелки, если строить линию пересечения с ними от точки $(-1,0)$, при дальшейшем повороте система снова становится неустойчивой.
\begin{figure}[H]
    \centering
    \includegraphics[width=0.6\linewidth]{lab_6/images3/nyquist_W4_tau_0.00.png}
    \caption{Годограф Найквиста $W_4(s)$}
    \label{5yn}
\end{figure}
\begin{figure}[H]
    \centering
    \includegraphics[width=0.6\linewidth]{lab_6/images3/nyquist_W4_tau_0.50.png}
    \caption{Годограф Найквиста $W_4(s)$}
    \label{6yn}
\end{figure}

Переходные характеристики для коэффициентов $\tau$, соответствующих различным областям значений коэффициента представлены на рис. \ref{t2}.
\\
Моделирование также подтверждает теорию и выводы, сделанные выше.
\begin{figure}[H]
    \centering
    \includegraphics[width=0.6\linewidth]{lab_6/images3/nyquist_W4_tau_0.70.png}
    \caption{Годограф Найквиста $W_4(s)$}
    \label{7yn}
\end{figure}

\begin{figure}[H]
    \centering
    \includegraphics[width=0.6\linewidth]{lab_6/images3/step_W4.png}
    \caption{Переходные характеристики $W_4(s)$}
    \label{t2}
\end{figure}


\section{Выводы}
В ходе лабораторной работы мы исследовали методы анализа устойчивости с использо-ванием критерия Найквиста. Мы выполненили три эксперимента, которые позволили проверить и подтвердить теоретические предпоположения, а также изучить влияние основных параметров на устойчивость замкнутой системы.

Мы проверили справедливость классического критерия Найквиста на трех объектах. Для каждого объекта были построены карты нулей и полюсов, переходные характеристики, годографы Найквиста и логарифмические частотные характеристики. Анализ показал полное соответствие между количеством оборотов годографа вокруг критической точки $(-1, 0)$ и разностью между количеством неустойчивых полюсов замкнутой и разомкнутой систем (формула $N = m - n$). Логарифмический критерий Найквиста, основанный на подсчете пересечений фазовочастотной характеристики критических отрезков при положительных значениях амплитудно-частотной характеристики, также подтвердил полученные результаты.

Мы исследовали влияние коэффициента усиления на устойчивость систем. Для двух различных передаточных функций с добавленным коэффициентом усиления $k$ были построены годографы Найквиста при различных значениях $k$ и определены области устойчивости. Проверено, что увеличение коэффициента усиления приводит к пропорци-ональному растяжению годографа, что вызывает изменение устойчивости при пересечении годографом критической точки $(-1, 0)$. Определены критические значения коэффициента усиления: для первой системы $W_1(s)$ устойчивость сохраняется при $k < 1$, для второй системы $W_2(s)$ — при $1 < k < 2$. Переходные характеристики, построенные для различных значений $k$, наглядно продемонстрировали изменение динамики системы от устойчивой к неустойчивой при выходе коэффициента усиления за пределы указанных областей.

Также мы проверили влияние чистого запаздывания на устойчивость систем. Добавле-ние звена запаздывания $e^{-\tau s}$ приводит к фазовому сдвигу, что проявляется в раскручивании годографа Найквиста. Для системы $W_3(s)$ определен критический коэффициент запаз-дывания $\tau \approx 0.254$, при превышении которого система теряет устойчивость. Для системы $W_4(s)$ обнаружена область устойчивости в интервале $0.218 < \tau < 0.598$, что соответствует случаю, когда критическая точка $(-1, 0)$ не охватывается годографом. Моделирование подтвердило предположения о границах устойчивости.

Данная работа показала практическую применимость критерия Найквиста, как инстру-мента анализа устойчивости систем, позволяющего не только оценивать устойчивость, но и определять запасы устойчивости, а также исследовать влияние различных параметров на динамические свойства систем.

\newpage
\section{Приложение}

\matlabtitle{Задание 1}
\begin{matlabcode}
% Система 1
R1 = [7 -9 52 51 35];
Q1 = [1 -5 20 -20 19 65];
D1 = [1 2 11 32 70 100];   
open_sys1 = tf(R1, Q1);
closed_sys1 = tf(R1, D1);
p1 = pole(open_sys1);
z1 = zero(open_sys1);
p_closed1 = pole(closed_sys1);
z_closed1 = zero(closed_sys1);

figure('Position', [100 100 1200 500]);
set(gcf, 'Color', 'w');

subplot(1,2,1);
hold on;
grid on;
xlabel('Re');
ylabel('Im');
title('Разомкнутая система 1', 'FontWeight', 'normal');

for i = 1:length(p1)
    if real(p1(i)) < 0
        plot(real(p1(i)), imag(p1(i)), 'xg', 'MarkerSize', 10, 'LineWidth', 2);
    else
        plot(real(p1(i)), imag(p1(i)), 'xr', 'MarkerSize', 10, 'LineWidth', 2);
    end
end
for i = 1:length(z1)
    plot(real(z1(i)), imag(z1(i)), 'bo', 'MarkerSize', 8, 'LineWidth', 2);
end
xlim([-2 3]);
ylim([-4 4]);
box on;
h1 = plot(NaN, NaN, 'xg', 'MarkerSize', 10, 'LineWidth', 2);
h2 = plot(NaN, NaN, 'xr', 'MarkerSize', 10, 'LineWidth', 2);
h3 = plot(NaN, NaN, 'bo', 'MarkerSize', 8, 'LineWidth', 2);
legend([h1, h2, h3], {'Уст. полюс', 'Неуст. полюс', 'Ноль'}, 'Location', 'best');

subplot(1,2,2);
hold on;
grid on;
xlabel('Re');
ylabel('Im');
title('Замкнутая система 1', 'FontWeight', 'normal');

for i = 1:length(p_closed1)
    if real(p_closed1(i)) < 0
        plot(real(p_closed1(i)), imag(p_closed1(i)), 'xg', 'MarkerSize', 10, 'LineWidth', 2);
    else
        plot(real(p_closed1(i)), imag(p_closed1(i)), 'xr', 'MarkerSize', 10, 'LineWidth', 2);
    end
end
for i = 1:length(z_closed1)
    plot(real(z_closed1(i)), imag(z_closed1(i)), 'bo', 'MarkerSize', 8, 'LineWidth', 2);
end
xlim([-3 2]);
ylim([-4 4]);
box on;
h1 = plot(NaN, NaN, 'xg', 'MarkerSize', 10, 'LineWidth', 2);
h2 = plot(NaN, NaN, 'xr', 'MarkerSize', 10, 'LineWidth', 2);
h3 = plot(NaN, NaN, 'bo', 'MarkerSize', 8, 'LineWidth', 2);
legend([h1, h2, h3], {'Уст. полюс', 'Неуст. полюс', 'Ноль'}, 'Location', 'best');

exportgraphics(gcf, "images/pz_combined1.png", Resolution=200);

figure('Position', [100 100 800 800]);
set(gcf, 'Color', 'w');

simtime = 15;
subplot(2,1,1);
[y_open1, t1] = step(open_sys1, simtime);
plot(t1, y_open1, 'b', 'LineWidth', 1.5);
xlabel('t, s');
ylabel('y');
grid on;
title('Разомкнутая система 1', 'FontWeight', 'normal');
xlim([0 simtime]);

subplot(2,1,2);
[y_closed1, t_closed1] = step(closed_sys1, simtime);
plot(t_closed1, y_closed1, 'b', 'LineWidth', 1.5);
xlabel('t, s');
ylabel('y');
grid on;
title('Замкнутая система 1', 'FontWeight', 'normal');
xlim([0 simtime]);

exportgraphics(gcf, "images/step_combined1.png", Resolution=200);

figure('Position', [100 100 800 600]);
set(gcf, 'Color', 'w');

w = logspace(-10, 10, 10000);
[re1, im1] = nyquist(open_sys1, w);
re1 = squeeze(re1);
im1 = squeeze(im1);

plot(re1, im1, 'b', 'LineWidth', 1.5);
hold on;
plot(re1, -im1, 'g', 'LineWidth', 1.5, 'LineStyle', '-');

arrow_spacing = 4000;
arrow_size = 0.1;

for k = 1:arrow_spacing:length(re1)-1
    x1 = re1(k);
    y1 = im1(k);
    x2 = re1(k+1);
    y2 = im1(k+1);

    theta = atan2(y2 - y1, x2 - x1);

    arrow_x = [0 -0.5 -0.5 0]*arrow_size;
    arrow_y = [0 0.2 -0.2 0]*arrow_size;

    R = [cos(theta) -sin(theta); sin(theta) cos(theta)];
    coords = R * [arrow_x; arrow_y];
    patch(x1 + coords(1,:), y1 + coords(2,:), 'r', 'EdgeColor','none');
end

plot(-1, 0, 'xr', 'MarkerSize', 9, 'LineWidth', 2);
xlabel('Re(W(j\omega))');
ylabel('Im(W(j\omega))');
grid on;
title('Годограф Найквиста 1', 'FontWeight', 'normal');
axis equal;
box on;
ax = gca;
ax.XGrid = 'on';
ax.YGrid = 'on';
ax.GridLineStyle = '-';
ax.GridAlpha = 0.3;
ax.LineWidth = 1.5;

exportgraphics(gcf, "images/nyquist1.png", Resolution=200);

figure('Position', [100 100 1000 800]);
set(gcf, 'Color', 'w');

[mag1, phase1, w1] = bode(open_sys1);
mag1 = squeeze(mag1);
phase1 = squeeze(phase1);

subplot(2,1,1);
semilogx(w1, 20*log10(mag1), 'b', 'LineWidth', 1.5);
grid on;
ylabel('L, dB');
title('ЛАЧХ разомкнутой системы 1', 'FontWeight', 'normal');

subplot(2,1,2);
semilogx(w1, phase1, 'r', 'LineWidth', 1.5);
grid on;
xlabel('\omega, rad/s');
ylabel('\phi, deg');
title('ЛФЧХ разомкнутой системы 1', 'FontWeight', 'normal');

exportgraphics(gcf, "images/bode1.png", Resolution=200);

% Система 2
R2 = [-9 -31 -73 121 432];
Q2 = [1 7 21 33 28 10];
D2 = [1 -2 -10 -40 149 442];
open_sys2 = tf(R2, Q2);
closed_sys2 = tf(R2, D2);
p2 = pole(open_sys2);
z2 = zero(open_sys2);
p_closed2 = pole(closed_sys2);
z_closed2 = zero(closed_sys2);


figure('Position', [100 100 1200 500]);
set(gcf, 'Color', 'w');

subplot(1,2,1);
hold on;
grid on;
xlabel('Re');
ylabel('Im');
title('Разомкнутая система 2', 'FontWeight', 'normal');

for i = 1:length(p2)
    if real(p2(i)) < 0
        plot(real(p2(i)), imag(p2(i)), 'xg', 'MarkerSize', 10, 'LineWidth', 2);
    else
        plot(real(p2(i)), imag(p2(i)), 'xr', 'MarkerSize', 10, 'LineWidth', 2);
    end
end
for i = 1:length(z2)
    plot(real(z2(i)), imag(z2(i)), 'bo', 'MarkerSize', 8, 'LineWidth', 2);
end
xlim([-3 3]);
ylim([-4 4]);
box on;
h1 = plot(NaN, NaN, 'xg', 'MarkerSize', 10, 'LineWidth', 2);
h2 = plot(NaN, NaN, 'xr', 'MarkerSize', 10, 'LineWidth', 2);
h3 = plot(NaN, NaN, 'bo', 'MarkerSize', 8, 'LineWidth', 2);
legend([h1, h2, h3], {'Уст. полюс', 'Неуст. полюс', 'Ноль'}, 'Location', 'best');

subplot(1,2,2);
hold on;
grid on;
xlabel('Re');
ylabel('Im');
title('Замкнутая система 2', 'FontWeight', 'normal');

for i = 1:length(p_closed2)
    if real(p_closed2(i)) < 0
        plot(real(p_closed2(i)), imag(p_closed2(i)), 'xg', 'MarkerSize', 10, 'LineWidth', 2);
    else
        plot(real(p_closed2(i)), imag(p_closed2(i)), 'xr', 'MarkerSize', 10, 'LineWidth', 2);
    end
end
for i = 1:length(z_closed2)
    plot(real(z_closed2(i)), imag(z_closed2(i)), 'bo', 'MarkerSize', 8, 'LineWidth', 2);
end
xlim([-3 5]);
ylim([-4 4]);
box on;
h1 = plot(NaN, NaN, 'xg', 'MarkerSize', 10, 'LineWidth', 2);
h2 = plot(NaN, NaN, 'xr', 'MarkerSize', 10, 'LineWidth', 2);
h3 = plot(NaN, NaN, 'bo', 'MarkerSize', 8, 'LineWidth', 2);
legend([h1, h2, h3], {'Уст. полюс', 'Неуст. полюс', 'Ноль'}, 'Location', 'best');

exportgraphics(gcf, "images/pz_combined2.png", Resolution=200);

figure('Position', [100 100 800 800]);
set(gcf, 'Color', 'w');

simtime = 20;
subplot(2,1,1);
[y_open2, t2] = step(open_sys2, simtime);
plot(t2, y_open2, 'b', 'LineWidth', 1.5);
xlabel('t, s');
ylabel('y');
grid on;
title('Разомкнутая система 2', 'FontWeight', 'normal');
xlim([0 simtime]);

subplot(2,1,2);
[y_closed2, t_closed2] = step(closed_sys2, simtime);
plot(t_closed2, y_closed2, 'b', 'LineWidth', 0.85);
xlabel('t, s');
ylabel('y');
grid on;
title('Замкнутая система 2', 'FontWeight', 'normal');
xlim([0 simtime]);

exportgraphics(gcf, "images/step_combined2.png", Resolution=200);

figure('Position', [100 100 800 600]);
set(gcf, 'Color', 'w');

w = logspace(-10, 10, 10000);
[re2, im2] = nyquist(open_sys2, w);
re2 = squeeze(re2);
im2 = squeeze(im2);

plot(re2, im2, 'b', 'LineWidth', 1.5);
hold on;
plot(re2, -im2, 'g', 'LineWidth', 1.5, 'LineStyle', '-');

arrow_spacing = 5000;
arrow_size = 5;

for k = 1:arrow_spacing:length(re2)-1
    x1 = re2(k);
    y1 = im2(k);
    x2 = re2(k+1);
    y2 = im2(k+1);

    theta = atan2(y2 - y1, x2 - x1);

    arrow_x = [0 -0.5 -0.5 0]*arrow_size;
    arrow_y = [0 0.2 -0.2 0]*arrow_size;

    R = [cos(theta) -sin(theta); sin(theta) cos(theta)];
    coords = R * [arrow_x; arrow_y];
    patch(x1 + coords(1,:), y1 + coords(2,:), 'r', 'EdgeColor','none');
end

plot(-1, 0, 'xr', 'MarkerSize', 9, 'LineWidth', 2);
xlabel('Re(W(j\omega))');
ylabel('Im(W(j\omega))');
grid on;
title('Годограф Найквиста 2', 'FontWeight', 'normal');
axis equal;
box on;
ax = gca;
ax.XGrid = 'on';
ax.YGrid = 'on';
ax.GridLineStyle = '-';
ax.GridAlpha = 0.3;
ax.LineWidth = 1.5;

exportgraphics(gcf, "images/nyquist2.png", Resolution=200);

figure('Position', [100 100 1000 800]);
set(gcf, 'Color', 'w');

[mag2, phase2, w2] = bode(open_sys2);
mag2 = squeeze(mag2);
phase2 = squeeze(phase2);

subplot(2,1,1);
semilogx(w2, 20*log10(mag2), 'b', 'LineWidth', 1.5);
grid on;
ylabel('L, dB');
title('ЛАЧХ разомкнутой системы 2', 'FontWeight', 'normal');

subplot(2,1,2);
semilogx(w2, phase2, 'r', 'LineWidth', 1.5);
grid on;
xlabel('\omega, rad/s');
ylabel('\phi, deg');
title('ЛФЧХ разомкнутой системы 2', 'FontWeight', 'normal');

exportgraphics(gcf, "images/bode2.png", Resolution=200);

% Система 3
R3 = [17 65 271 528 390];
Q3 = [1 -5 9 -3 -8 10];
D3 = [1 12 74 268 520 400];
open_sys3 = tf(R3, Q3);
closed_sys3 = tf(R3, D3);
p3 = pole(open_sys3);
z3 = zero(open_sys3);
p_closed3 = pole(closed_sys3);
z_closed3 = zero(closed_sys3);



figure('Position', [100 100 1200 500]);
set(gcf, 'Color', 'w');

subplot(1,2,1);
hold on;
grid on;
xlabel('Re');
ylabel('Im');
title('Разомкнутая система 3', 'FontWeight', 'normal');

for i = 1:length(p3)
    if real(p3(i)) < 0
        plot(real(p3(i)), imag(p3(i)), 'xg', 'MarkerSize', 10, 'LineWidth', 2);
    else
        plot(real(p3(i)), imag(p3(i)), 'xr', 'MarkerSize', 10, 'LineWidth', 2);
    end
end
for i = 1:length(z3)
    plot(real(z3(i)), imag(z3(i)), 'bo', 'MarkerSize', 8, 'LineWidth', 2);
end
xlim([-2 3]);
ylim([-4 4]);
box on;
h1 = plot(NaN, NaN, 'xg', 'MarkerSize', 10, 'LineWidth', 2);
h2 = plot(NaN, NaN, 'xr', 'MarkerSize', 10, 'LineWidth', 2);
h3 = plot(NaN, NaN, 'bo', 'MarkerSize', 8, 'LineWidth', 2);
legend([h1, h2, h3], {'Уст. полюс', 'Неуст. полюс', 'Ноль'}, 'Location', 'best');

subplot(1,2,2);
hold on;
grid on;
xlabel('Re');
ylabel('Im');
title('Замкнутая система 3', 'FontWeight', 'normal');

for i = 1:length(p_closed3)
    if real(p_closed3(i)) < 0
        plot(real(p_closed3(i)), imag(p_closed3(i)), 'xg', 'MarkerSize', 10, 'LineWidth', 2);
    else
        plot(real(p_closed3(i)), imag(p_closed3(i)), 'xr', 'MarkerSize', 10, 'LineWidth', 2);
    end
end
for i = 1:length(z_closed3)
    plot(real(z_closed3(i)), imag(z_closed3(i)), 'bo', 'MarkerSize', 8, 'LineWidth', 2);
end
xlim([-4 0]);
ylim([-5 5]);
box on;
h1 = plot(NaN, NaN, 'xg', 'MarkerSize', 10, 'LineWidth', 2);
h2 = plot(NaN, NaN, 'xr', 'MarkerSize', 10, 'LineWidth', 2);
h3 = plot(NaN, NaN, 'bo', 'MarkerSize', 8, 'LineWidth', 2);
legend([h1, h2, h3], {'Уст. полюс', 'Неуст. полюс', 'Ноль'}, 'Location', 'best');

exportgraphics(gcf, "images/pz_combined3.png", Resolution=200);

figure('Position', [100 100 800 800]);
set(gcf, 'Color', 'w');

simtime = 10;
subplot(2,1,1);
[y_open3, t3] = step(open_sys3, simtime);
plot(t3, y_open3, 'b', 'LineWidth', 1.5);
xlabel('t, s');
ylabel('y');
grid on;
title('Разомкнутая система 3', 'FontWeight', 'normal');
xlim([0 simtime]);

subplot(2,1,2);
[y_closed3, t_closed3] = step(closed_sys3, simtime);
plot(t_closed3, y_closed3, 'b', 'LineWidth', 1.5);
xlabel('t, s');
ylabel('y');
grid on;
title('Замкнутая система 3', 'FontWeight', 'normal');
xlim([0 simtime]);

exportgraphics(gcf, "images/step_combined3.png", Resolution=200);

figure('Position', [100 100 800 600]);
set(gcf, 'Color', 'w');

w = logspace(-10, 10, 10000);
[re3, im3] = nyquist(open_sys3, w);
re3 = squeeze(re3);
im3 = squeeze(im3);

plot(re3, im3, 'b', 'LineWidth', 1.5);
hold on;
plot(re3, -im3, 'g', 'LineWidth', 1.5, 'LineStyle', '-');

arrow_spacing = 5000;
arrow_size = 5;

for k = 1:arrow_spacing:length(re3)-1
    x1 = re3(k);
    y1 = im3(k);
    x2 = re3(k+1);
    y2 = im3(k+1);

    theta = atan2(y2 - y1, x2 - x1);

    arrow_x = [0 -0.5 -0.5 0]*arrow_size;
    arrow_y = [0 0.2 -0.2 0]*arrow_size;

    R = [cos(theta) -sin(theta); sin(theta) cos(theta)];
    coords = R * [arrow_x; arrow_y];
    patch(x1 + coords(1,:), y1 + coords(2,:), 'r', 'EdgeColor','none');
end

plot(-1, 0, 'xr', 'MarkerSize', 9, 'LineWidth', 2);
xlabel('Re(W(j\omega))');
ylabel('Im(W(j\omega))');
grid on;
title('Годограф Найквиста 3', 'FontWeight', 'normal');
axis equal;
box on;
ax = gca;
ax.XGrid = 'on';
ax.YGrid = 'on';
ax.GridLineStyle = '-';
ax.GridAlpha = 0.3;
ax.LineWidth = 1.5;

exportgraphics(gcf, "images/nyquist3.png", Resolution=200);

figure('Position', [100 100 1000 800]);
set(gcf, 'Color', 'w');

[mag3, phase3, w3] = bode(open_sys3);
mag3 = squeeze(mag3);
phase3 = squeeze(phase3);

subplot(2,1,1);
semilogx(w3, 20*log10(mag3), 'b', 'LineWidth', 1.5);
grid on;
ylabel('L, dB');
title('ЛАЧХ разомкнутой системы 3', 'FontWeight', 'normal');

subplot(2,1,2);
semilogx(w3, phase3, 'r', 'LineWidth', 1.5);
grid on;
xlabel('\omega, rad/s');
ylabel('\phi, deg');
title('ЛФЧХ разомкнутой системы 3', 'FontWeight', 'normal');

exportgraphics(gcf, "images/bode3.png", Resolution=200);
\end{matlabcode}

\matlabtitle{Задание 2}
\begin{matlabcode}
R1 = [1 -1];
Q1 = [1 3 1];
k_values = [0.7 1 1.5 3];

for idx = 1:length(k_values)
    k = k_values(idx);
    
    open_sys = tf(k * R1, Q1);
    closed_sys = feedback(open_sys, 1);
    
    w = logspace(-10, 10, 10000);
    [re, im] = nyquist(open_sys, w);
    re = squeeze(re);
    im = squeeze(im);
    
    x_range = max(re) - min(re);
    arrow_size = 0.1 * x_range^0.7;
    
    figure('Position', [100 100 800 600], 'Visible', 'off');
    set(gcf, 'Color', 'w');
    
    plot(re, im, 'b', 'LineWidth', 1.5);
    hold on;
    plot(re, -im, 'g', 'LineWidth', 1.5, 'LineStyle', '-');
    
    arrow_spacing = 5000;
    
    for k_arr = 1:arrow_spacing:length(re)-1
        x1 = re(k_arr);
        y1 = im(k_arr);
        x2 = re(k_arr+1);
        y2 = im(k_arr+1);
        
        theta = atan2(y2 - y1, x2 - x1);
        arrow_x = [0 -0.5 -0.5 0]*arrow_size;
        arrow_y = [0 0.2 -0.2 0]*arrow_size;
        
        R = [cos(theta) -sin(theta); sin(theta) cos(theta)];
        coords = R * [arrow_x; arrow_y];
        patch(x1 + coords(1,:), y1 + coords(2,:), 'r', 'EdgeColor','none');
    end
    
    plot(-1, 0, 'xr', 'MarkerSize', 9, 'LineWidth', 2);
    xlabel('Re(W(j\omega))');
    ylabel('Im(W(j\omega))');
    grid on;
    axis equal;
    box on;
    title(['W1(s) годограф Найквиста, k = ' num2str(k)], 'FontWeight', 'normal');
    
    exportgraphics(gcf, sprintf("images2/nyquist_W1_k_%.1f.png", k), Resolution=200);
    close(gcf);
end

figure('Position', [100 100 800 1000], 'Visible', 'off');
set(gcf, 'Color', 'w');

simtime = 150;
for idx = 1:length(k_values)
    k = k_values(idx);
    open_sys = tf(k * R1, Q1);
    closed_sys = feedback(open_sys, 1);
    [y_closed, t_closed] = step(closed_sys, simtime);
    
    subplot(length(k_values), 1, idx);
    plot(t_closed, y_closed, 'r', 'LineWidth', 2);
    
    grid on;
    xlabel('t, s');
    ylabel('y');
    title(['W1(s), k = ' num2str(k)], 'FontWeight', 'normal');
    xlim([0 simtime]);
end

exportgraphics(gcf, "images2/step_W1.png", Resolution=200);
close(gcf);

R2 = [10 -3 13 -2];
Q2 = [10 8 -5 4];
k_values = [0.5 1 1.5 2 3];

for idx = 1:length(k_values)
    k = k_values(idx);
    
    open_sys = tf(k * R2, Q2);
    closed_sys = feedback(open_sys, 1);
    
    w = logspace(-10, 10, 10000);
    [re, im] = nyquist(open_sys, w);
    re = squeeze(re);
    im = squeeze(im);
    
    x_range = max(re) - min(re);
    arrow_size = 0.1 * x_range^0.7;
    
    figure('Position', [100 100 800 600], 'Visible', 'off');
    set(gcf, 'Color', 'w');
    
    plot(re, im, 'b', 'LineWidth', 1.5);
    hold on;
    plot(re, -im, 'g', 'LineWidth', 1.5, 'LineStyle', '-');
    
    arrow_spacing = 5000;
    
    for k_arr = 1:arrow_spacing:length(re)-1
        x1 = re(k_arr);
        y1 = im(k_arr);
        x2 = re(k_arr+1);
        y2 = im(k_arr+1);
        
        theta = atan2(y2 - y1, x2 - x1);
        arrow_x = [0 -0.5 -0.5 0]*arrow_size;
        arrow_y = [0 0.2 -0.2 0]*arrow_size;
        
        R = [cos(theta) -sin(theta); sin(theta) cos(theta)];
        coords = R * [arrow_x; arrow_y];
        patch(x1 + coords(1,:), y1 + coords(2,:), 'r', 'EdgeColor','none');
    end
    
    plot(-1, 0, 'xr', 'MarkerSize', 9, 'LineWidth', 2);
    xlabel('Re(W(j\omega))');
    ylabel('Im(W(j\omega))');
    grid on;
    axis equal;
    box on;
    title(['W2(s) годограф Найквиста, k = ' num2str(k)], 'FontWeight', 'normal');
    
    exportgraphics(gcf, sprintf("images2/nyquist_W2_k_%.1f.png", k), Resolution=200);
    close(gcf);
end

figure('Position', [100 100 800 1000], 'Visible', 'off');
set(gcf, 'Color', 'w');

simtime = 150;
for idx = 1:length(k_values)
    k = k_values(idx);
    open_sys = tf(k * R2, Q2);
    closed_sys = feedback(open_sys, 1);
    [y_closed, t_closed] = step(closed_sys, simtime);
    
    subplot(length(k_values), 1, idx);
    plot(t_closed, y_closed, 'r', 'LineWidth', 2);
    
    grid on;
    xlabel('t, s');
    ylabel('y');
    title(['W2(s), k = ' num2str(k)], 'FontWeight', 'normal');
    xlim([0 simtime]);
end

exportgraphics(gcf, "images2/step_W2.png", Resolution=200);
close(gcf);

\end{matlabcode}

\matlabtitle{Задание 3}
\begin{matlabcode}
R3 = [2 9];
Q3 = [1 1 9];
tau_values = [0.0 0.25 0.5];

for idx = 1:length(tau_values)
    tau = tau_values(idx);
    
    open_sys = tf(R3, Q3, 'InputDelay', tau);
    closed_sys = feedback(open_sys, 1);
    
    w = logspace(-10, 10, 10000);
    [re, im] = nyquist(open_sys, w);
    re = squeeze(re);
    im = squeeze(im);
    
    x_range = max(re) - min(re);
    arrow_size = 0.1 * x_range^0.7;
    
    figure('Position', [100 100 800 600], 'Visible', 'off');
    set(gcf, 'Color', 'w');
    
    plot(re, im, 'b', 'LineWidth', 1.5);
    hold on;
    plot(re, -im, 'g', 'LineWidth', 1.5, 'LineStyle', '-');
    
    arrow_spacing = 5000;
    
    for k_arr = 1:arrow_spacing:length(re)-1
        x1 = re(k_arr);
        y1 = im(k_arr);
        x2 = re(k_arr+1);
        y2 = im(k_arr+1);
        
        theta = atan2(y2 - y1, x2 - x1);
        arrow_x = [0 -0.5 -0.5 0]*arrow_size;
        arrow_y = [0 0.2 -0.2 0]*arrow_size;
        
        R = [cos(theta) -sin(theta); sin(theta) cos(theta)];
        coords = R * [arrow_x; arrow_y];
        patch(x1 + coords(1,:), y1 + coords(2,:), 'r', 'EdgeColor','none');
    end
    
    plot(-1, 0, 'xr', 'MarkerSize', 9, 'LineWidth', 2);
    xlabel('Re(W(j\omega))');
    ylabel('Im(W(j\omega))');
    grid on;
    axis equal;
    box on;
    title(['W3(s) годограф Найквиста, tau = ' num2str(tau)], 'FontWeight', 'normal');
    
    exportgraphics(gcf, sprintf("images3/nyquist_W3_tau_%.2f.png", tau), Resolution=200);
    close(gcf);
end

figure('Position', [100 100 800 1000], 'Visible', 'off');
set(gcf, 'Color', 'w');

simtime = 150;
for idx = 1:length(tau_values)
    tau = tau_values(idx);
    open_sys = tf(R3, Q3, 'InputDelay', tau);
    closed_sys = feedback(open_sys, 1);
    [y_closed, t_closed] = step(closed_sys, simtime);
    
    subplot(length(tau_values), 1, idx);
    plot(t_closed, y_closed, 'r', 'LineWidth', 2);
    
    grid on;
    xlabel('t, s');
    ylabel('y');
    title(['W3(s), tau = ' num2str(tau)], 'FontWeight', 'normal');
    xlim([0 simtime]);
end

exportgraphics(gcf, "images3/step_W3.png", Resolution=200);
close(gcf);

R4 = [10 -10 13];
Q4 = [10 0 37 25];
tau_values = [0.0 0.5 0.7];

for idx = 1:length(tau_values)
    tau = tau_values(idx);
    
    open_sys = tf(R4, Q4, 'InputDelay', tau);
    closed_sys = feedback(open_sys, 1);
    
    w = logspace(-10, 10, 10000);
    [re, im] = nyquist(open_sys, w);
    re = squeeze(re);
    im = squeeze(im);
    
    x_range = max(re) - min(re);
    arrow_size = 0.1 * x_range^0.7;
    
    figure('Position', [100 100 800 600], 'Visible', 'off');
    set(gcf, 'Color', 'w');
    
    plot(re, im, 'b', 'LineWidth', 1.5);
    hold on;
    plot(re, -im, 'g', 'LineWidth', 1.5, 'LineStyle', '-');
    
    arrow_spacing = 5000;
    
    for k_arr = 1:arrow_spacing:length(re)-1
        x1 = re(k_arr);
        y1 = im(k_arr);
        x2 = re(k_arr+1);
        y2 = im(k_arr+1);
        
        theta = atan2(y2 - y1, x2 - x1);
        arrow_x = [0 -0.5 -0.5 0]*arrow_size;
        arrow_y = [0 0.2 -0.2 0]*arrow_size;
        
        R = [cos(theta) -sin(theta); sin(theta) cos(theta)];
        coords = R * [arrow_x; arrow_y];
        patch(x1 + coords(1,:), y1 + coords(2,:), 'r', 'EdgeColor','none');
    end
    
    plot(-1, 0, 'xr', 'MarkerSize', 9, 'LineWidth', 2);
    xlabel('Re(W(j\omega))');
    ylabel('Im(W(j\omega))');
    grid on;
    axis equal;
    box on;
    title(['W4(s) годограф Найквиста, tau = ' num2str(tau)], 'FontWeight', 'normal');
    
    exportgraphics(gcf, sprintf("images3/nyquist_W4_tau_%.2f.png", tau), Resolution=200);
    close(gcf);
end

figure('Position', [100 100 800 1000], 'Visible', 'off');
set(gcf, 'Color', 'w');

simtime = 150;
for idx = 1:length(tau_values)
    tau = tau_values(idx);
    open_sys = tf(R4, Q4, 'InputDelay', tau);
    closed_sys = feedback(open_sys, 1);
    [y_closed, t_closed] = step(closed_sys, simtime);
    
    subplot(length(tau_values), 1, idx);
    plot(t_closed, y_closed, 'r', 'LineWidth', 2);
    
    grid on;
    xlabel('t, s');
    ylabel('y');
    title(['W4(s), tau = ' num2str(tau)], 'FontWeight', 'normal');
    xlim([0 simtime]);
end

exportgraphics(gcf, "images3/step_W4.png", Resolution=200);
close(gcf);

\end{matlabcode}
